(2)
Assume that $\mathfrak g_1\ne\{0\}$ and $\mu\geqq3$.
By transitivity of $\mathfrak g$,
$[\mathfrak g_1,\mathfrak e]\ne\{0\}$ or $[\mathfrak g_1,\mathfrak f]\ne\{0\}$.
We may assume that $[\mathfrak g_1,\mathfrak e]\ne\{0\}$.
Then there exists an irreducible component $\mathfrak g'_1$ of
the $\mathfrak g_0$-module~$\mathfrak g_1$
such that $[\mathfrak g'_1,\mathfrak e]\ne\{0\}$ and
$[\mathfrak g'_1,\mathfrak f]=\{0\}$.
The subalgebra $\mathfrak e\oplus[\mathfrak e,\mathfrak g'_1]\oplus
\mathfrak g'_1$ is a simple GLA of the f\/irst kind.
Since $\mathfrak g_0$ is isomorphic to $\mathfrak{gl}(\mathfrak e)\oplus
\mathfrak{gl}(\mathfrak f)$,
$\mathfrak e\oplus[\mathfrak e,\mathfrak g'_1]\oplus \mathfrak g'_1$ is of type $(A_m,\{\alpha_1\})$.
Let~$D$ be a nonzero element of $\mathfrak g'_1$.
There exist $\lambda\in \mathfrak e^*$ and $\eta\in\mathfrak f^*$ such that
\[
[[D,Z],U]=\lambda(U)Z+\lambda(Z)U, \qquad [[D,Z],W]=\eta(Z)W,
\]
where $Z,U\in \mathfrak e$ and $W\in \mathfrak f$
(cf.~\cite[p.~4]{tan57:1}).
Let $X$ (resp.~$Y$) be a nonzero element of $\mathfrak e$
(resp. $\mathfrak f$).
Since the subalgebra generated by $X,Y$ is a free FGLA of type $(2,\mu)$
(Remark~\ref{rem8.1}~(2)),
\begin{gather*}
\ad(X)^\mu(Y)=0, \qquad \ad(X)^{\mu-1}(Y)\ne 0, \\
\ad(Y)\ad(X)^{\mu-1}(Y)=0, \qquad \ad(Y)\ad(X)^{\mu-2}(Y)\ne 0
\end{gather*}
(Lemma \ref{lem2.1}).
By induction on $\mu$, we see that
\begin{gather*}
0  = \ad(D)\ad(X)^\mu(Y)=(\mu(\mu-1)\lambda(X)+\mu\eta(X))\ad(X)^{\mu-1}(Y),
\\
0  = \ad(D)\ad(Y)\ad(X)^{\mu-1}(Y)\\
\hphantom{0}{} =((\mu-1)(\mu-2)\lambda(X)
+(\mu-1)\eta(X))\ad(Y)\ad(X)^{\mu-2}(Y).
\end{gather*}
Since
\[
\det\begin{bmatrix}
\mu(\mu-1) & \mu \\
(\mu-1)(\mu-2) & \mu-1
\end{bmatrix}=
\mu(\mu-1)\ne0,
\]
we see that $\lambda(X)=\eta(X)=0$, which is a contradiction.
Thus we obtain that $\mu=2$ if \mbox{$\dim \mathfrak g_1\ne\{0\}$}.
From Example~\ref{ex8.2}, it follows that
$\gla g$ is a simple GLA of type $(A_{m+n}$, $\{\alpha_m,\alpha_{m+1}\})$
if $\dim \mathfrak g_1\ne\{0\}$.
\end{proof}
